
Reinforcement learning from execution feedback for code generation suffers from sparse rewards: binary test pass/fail signals provide no gradient for partial progress, and all tokens receive uniform credit regardless of their contribution to the outcome. Fine-Grained Optimization (FGO) addresses this by masking non-executed tokens from gradient updates, but prior work conflates FGO with other components (curriculum learning, feedback content), making it impossible to isolate the mechanism's contribution.

This paper presents a controlled mechanism validation study of FGO. The technique is decomposed into three testable components: (1) execution trace collection, (2) gradient exclusion via token masking, and (3) credit assignment to executed tokens. For each component, explicit falsification criteria are established and the mechanism is verified empirically.

Experiments on HumanEval and MBPP demonstrate that: trace collection achieves 100\% capture rate using Python's \texttt{sys.settrace}; gradient exclusion is correct, with masked tokens receiving exactly zero gradient in all verification checks; and FGO improves final pass@1 by approximately 10\% (0.244 vs. 0.222) in simulation-based evaluation. Token-level classification achieves 81\% F1, below the 95\% target, due to tokenizer-to-line mapping misalignment rather than trace collection failure.

The central finding is that FGO's benefit derives from gradient exclusion: non-executed tokens receive exactly zero gradient, while executed tokens receive 1.78x signal concentration. This mechanism validation approach---decomposing techniques into independently testable components with falsification criteria---provides a template for rigorous evaluation of credit assignment methods in code RL. Claims regarding convergence speed and content-versus-granularity comparisons remain unverified and are deferred to future work requiring full training runs.
